% ============================================================
% results.tex — Sección IX: Resultados
% ============================================================

\section{Resultados}
\label{sec:results}

Esta sección presenta los resultados experimentales obtenidos en la evaluación de \sirca{} frente a los baselines definidos, organizados según las tres dimensiones de evaluación: recuperación semántica, fidelidad generativa y reducción de alucinaciones.

\subsection{Resultados de recuperación semántica}
\label{subsec:res_retrieval}

La Tabla~\ref{tab:retrieval_results} presenta los resultados de las métricas de recuperación para cada sistema evaluado, promediados sobre las 60 consultas del conjunto de evaluación.

\begin{table}[!t]
\centering
\caption{Resultados de recuperación semántica (promedio $\pm$ desviación estándar sobre 60 consultas).}
\label{tab:retrieval_results}
\renewcommand{\arraystretch}{1.3}
\footnotesize
\begin{tabular}{l c c c}
\toprule
\textbf{Sistema} & \textbf{Recall@5} & \textbf{P@5} & \textbf{MRR} \\
\midrule
B1: BM25 puro             & 0.42 $\pm$ 0.11 & 0.38 $\pm$ 0.09 & 0.51 $\pm$ 0.14 \\
B2: RAG genérico           & 0.58 $\pm$ 0.10 & 0.52 $\pm$ 0.08 & 0.64 $\pm$ 0.12 \\
B3: RAG + BioBERT          & 0.67 $\pm$ 0.09 & 0.61 $\pm$ 0.08 & 0.73 $\pm$ 0.11 \\
\midrule
\textbf{\sirca{}}          & \textbf{0.78 $\pm$ 0.07} & \textbf{0.71 $\pm$ 0.06} & \textbf{0.84 $\pm$ 0.08} \\
\bottomrule
\end{tabular}
\end{table}

Los resultados muestran que \sirca{} supera consistentemente a todos los baselines en las tres métricas de recuperación. Comparado con BM25 puro (B1), \sirca{} obtiene una mejora relativa del 85.7\% en Recall@5, lo cual evidencia la superioridad del retrieval híbrido con embeddings biomédicos frente a la recuperación léxica. Respecto al RAG genérico (B2), la mejora del 34.5\% en Recall@5 demuestra el beneficio de la especialización biomédica en los embeddings y el reranking contextual. Frente al RAG con BioBERT (B3), la mejora del 16.4\% en Recall@5 se atribuye a la combinación de retrieval híbrido y reranking contextual, que captura tanto coincidencias terminológicas exactas como relaciones semánticas profundas.

El análisis por tipo de consulta revela que las mejoras más significativas se observan en consultas exploratorias (mejora del 22.1\% en Recall@5 frente a B3), donde la recuperación semántica es particularmente ventajosa para identificar documentos relevantes que utilizan terminología variada. Las consultas factuales muestran mejoras moderadas (11.3\%), ya que el componente BM25 del retrieval híbrido contribuye eficazmente a la localización de términos específicos.

\subsection{Resultados de fidelidad generativa}
\label{subsec:res_faithfulness}

La Tabla~\ref{tab:generation_results} presenta las métricas de fidelidad generativa, comparando las respuestas generadas por cada sistema frente a las respuestas de referencia elaboradas por expertos.

\begin{table}[!t]
\centering
\caption{Resultados de fidelidad generativa (promedio sobre 60 consultas).}
\label{tab:generation_results}
\renewcommand{\arraystretch}{1.3}
\footnotesize
\begin{tabular}{l c c c}
\toprule
\textbf{Sistema} & \textbf{BERTScore} & \textbf{ROUGE-L} & \textbf{Sim. Sem.} \\
\midrule
B2: RAG genérico   & 0.71 & 0.34 & 0.68 \\
B3: RAG + BioBERT  & 0.76 & 0.39 & 0.73 \\
B4: LLM directo    & 0.62 & 0.22 & 0.57 \\
\midrule
\textbf{\sirca{}}  & \textbf{0.82} & \textbf{0.46} & \textbf{0.79} \\
\bottomrule
\end{tabular}
\end{table}

\sirca{} obtiene los valores más altos en todas las métricas de fidelidad generativa. La mejora en BERTScore (0.82 vs.\ 0.76 de B3) indica que las respuestas generadas por \sirca{} capturan con mayor fidelidad la semántica de las respuestas de referencia. La mejora sustancial en ROUGE-L (0.46 vs.\ 0.39) sugiere que la incorporación de documentos más relevantes mediante retrieval híbrido y reranking produce respuestas con mayor solapamiento léxico con las referencias.

El resultado del LLM directo (B4) confirma las limitaciones de los modelos generativos cuando operan sin acceso a evidencia documental externa. El BERTScore de 0.62 y la similitud semántica de 0.57 reflejan la tendencia del modelo a generar respuestas genéricas que carecen de la especificidad requerida por consultas sobre farmacología de plantas medicinales peruanas.

\subsection{Resultados de grounding y reducción de alucinaciones}
\label{subsec:res_hallucination}

La Tabla~\ref{tab:grounding_results} presenta las métricas de grounding documental y tasas de alucinación.

\begin{table}[!t]
\centering
\caption{Resultados de grounding documental y alucinación (evaluación sobre 60 consultas).}
\label{tab:grounding_results}
\renewcommand{\arraystretch}{1.3}
\footnotesize
\begin{tabular}{l c c}
\toprule
\textbf{Sistema} & \textbf{Tasa de Grounding} & \textbf{Tasa de Alucinación} \\
\midrule
B2: RAG genérico   & 0.61 & 0.23 \\
B3: RAG + BioBERT  & 0.69 & 0.18 \\
B4: LLM directo    & 0.12 & 0.54 \\
\midrule
\textbf{\sirca{}}  & \textbf{0.84} & \textbf{0.09} \\
\bottomrule
\end{tabular}
\end{table}

\sirca{} alcanza una tasa de grounding de 0.84, indicando que el 84\% de las afirmaciones factuales en las respuestas generadas están fundamentadas en los documentos recuperados. Simultáneamente, la tasa de alucinación se reduce a 0.09, lo cual representa una mejora del 83.3\% respecto al LLM directo (0.54) y del 50.0\% respecto al RAG con BioBERT (0.18).

Estos resultados confirman que la combinación de retrieval contextual de alta calidad, instrucciones de grounding explícitas y verificación posterior de fundamentación documental es efectiva para la reducción de alucinaciones en el dominio biomédico especializado.

\subsection{Análisis de latencia}
\label{subsec:res_latency}

La Tabla~\ref{tab:latency} presenta los tiempos de respuesta promedio para cada etapa del pipeline de \sirca{}.

\begin{table}[!t]
\centering
\caption{Latencia promedio por etapa del pipeline (ms).}
\label{tab:latency}
\renewcommand{\arraystretch}{1.2}
\footnotesize
\begin{tabular}{l r}
\toprule
\textbf{Etapa} & \textbf{Latencia (ms)} \\
\midrule
Embedding de consulta        & 45 \\
Recuperación híbrida (K=20)  & 120 \\
Reranking contextual (K'=5)  & 350 \\
Generación LLM               & 2\,800 \\
\midrule
\textbf{Total end-to-end}    & \textbf{3\,315} \\
\bottomrule
\end{tabular}
\end{table}

El tiempo total end-to-end de 3.3 segundos resulta aceptable para un sistema de consulta interactivo, con la generación del LLM constituyendo el cuello de botella principal (84.5\% del tiempo total). La etapa de recuperación híbrida y reranking añade un overhead moderado (470~ms) que se justifica por la mejora significativa en calidad de recuperación.
